\begin{lemma}
\label{lemma-subalgebra-geometrically-irreducible}
Let $k$ be a field. Let $S$ be a $k$-algebra.
\begin{enumerate}
\item If $S$ is geometrically irreducible over $k$ so is every
$k$-subalgebra.
\item If all finitely generated $k$-subalgebras of $S$ are
geometrically irreducible, then $S$ is geometrically irreducible.
\item A directed colimit of geometrically irreducible $k$-algebras
is geometrically irreducible.
\end{enumerate}
\end{lemma}

\begin{proof}
For a nonzero ring $A$, its spectrum
is irreducible exactly when its
nilradical is prime, by Lemma
\ref{lemma-irreducible}. Equivalently,
whenever $ab$ is nilpotent, at
least one of $a,b$ is nilpotent.
We keep nonzeroness in this
criterion; the zero ring has
empty spectrum and is not
irreducible.

For (1), retain the original
subalgebra $S'\subseteq S$ and
any original field extension
$L/k$. The map
$S'\otimes_k L\to S\otimes_k L$
is injective by extension of
a vector-space basis. For any
injective ring map $A\to B$,
the nilradical of $A$ is exactly
the inverse image of that of
$B$: a power becomes zero in
$B$ exactly when it was zero
in $A$. Here the receiving
nilradical is a proper prime,
so its inverse image is also
a proper prime. The source is
nonzero because its $1$ remains
the nonzero $1$ of the target.
The criterion proves that
$S'\otimes_k L$ is irreducible,
for every such $L$, proving (1).
This argument works for any
injective $k$-algebra map, with
its actual image subalgebra.

For (3), let $(S_i)$ be a
nonempty directed system of
geometrically irreducible
$k$-algebras with their original
unital transition maps; they
need not be injective. Put
$S=\mathop{\rm colim}_i S_i$.
For each original $L/k$ there
are inverse maps
$$
S\otimes_k L\longrightarrow
\mathop{\rm colim}_i(S_i\otimes_k L),
\qquad [i,s]\otimes a\longmapsto[i,s\otimes a],
$$
and
$$
[i,\sum_j s_j\otimes a_j]
\longmapsto\sum_j[i,s_j]\otimes a_j.
$$
All expressions use finitely
many elements, so a common
later index verifies balancing,
independence of representatives,
and both composites. Let $A_i$
denote the original $S_i\otimes_k L$.
Each is nonzero with prime
nilradical. Their unital colimit
$A$ is nonzero: an equality
$1=0$ there would hold at a
later stage, contradicting
$1\neq0$ in that $A_i$.

Suppose original elements
$a,b\in A$ satisfy $(ab)^n=0$.
Choose representatives at a
common stage $i$. The zero
equality holds at some later
stage $j$. Its nilradical is
prime, so the image of $a$
or of $b$ is nilpotent at
that stage. The corresponding
element in the original
colimit is then nilpotent.
This proves that the nilradical
of $A$ is prime, and proves
(3) after the explicit tensor
comparison. In particular no
injection assumption on the
transition maps is used.
If an empty colimit is allowed
in the category of $k$-algebras,
it is its initial object $k$,
which also satisfies the
assertion.

For (2), the original finitely
generated $k$-subalgebras of $S$
form a directed system under
inclusion: adjoining the union
of two finite lists gives a
common member. Every original
element of $S$ occurs in such
a subalgebra, so their colimit
is exactly $S$. Apply (3).
The hypothesis excludes $S=0$,
since its finitely generated
subalgebra generated by the
image of $k$ would itself be
zero and would not satisfy
the stated hypothesis.

The same nilradical criterion
also proves localization
stability with its necessary
nonemptiness qualification.
Assume $S$ is geometrically irreducible.
Let $W\subseteq S$ be a
multiplicative subset and
assume $W^{-1}S\neq0$.
Since $S$ is geometrically
irreducible, its nilradical
$N$ is prime; $W$ avoids $N$,
because a nilpotent becoming
invertible would make the
localization zero. For every
original $L/k$, the injection
$S\to S\otimes_k L$ shows
that no image of an element
of $W$ is nilpotent. The
prime nilradical of that
tensor therefore survives
localization and gives a
nonzero ring with one minimal
prime. Explicitly the nilradical
of a localization $B_W$ is
$W^{-1}\operatorname{Nil}(B)$:
if $(b/w)^n=0$, some $v\in W$
has $v b^n=0$, and then
$(vb)^n=v^{n-1}(v b^n)=0$,
while $b/w=(vb)/(vw)$.
The converse follows by taking
powers of a nilpotent numerator.
The original tensor-localization
maps are
$$
(W^{-1}S)\otimes_k L
\longrightarrow (W\otimes1)^{-1}(S\otimes_k L),
\qquad(s/w)\otimes a\longmapsto(s\otimes a)/(w\otimes1),
$$
and the inverse sends
$\sum_i(s_i\otimes a_i)/(w\otimes1)$
to $\sum_i(s_i/w)\otimes a_i$.
The two localization universal
properties prove these maps
well-defined and inverse.
Thus every nonzero localization
of the original $S$ remains
geometrically irreducible.
\end{proof}